\subsection{§II.4 --- Assainissement : réparation de \texttt{familles\_algebre} (import aliasé)}
\textbf{Énoncé (Bourbaki).} Aucun --- il ne s'agit pas d'un résultat mathématique mais d'une opération
d'hygiène du dépôt.

\textbf{Formalisé.} Le module \texttt{ensembles\_familles\_algebre} réutilise l'infrastructure
famille-réciproque déjà certifiée, exposée sous ses noms canoniques \texttt{membre\_image\_recip} et
\texttt{famille\_image\_recip}. L'import a été remis en cohérence par aliasage
(\texttt{membre\_image\_recip as membre\_image\_reciproque}, \texttt{famille\_image\_recip as famille\_reciproque}),
ce qui a supprimé la dernière erreur de collecte du dépôt sous \texttt{pytest}. \\
Module : \texttt{bourbaki/ensembles/familles/ii\_4\_reunion\_intersection\_familles/ii\_4\_1\_definitions\_algebre/ensembles\_familles\_algebre.py}.

\textbf{Preuve (\Nstar).} Sans objet : aucun théorème n'est construit ni modifié par cette opération.
Les symboles importés ne sont que des renommages locaux des fonctions certifiées d'origine ; la chaîne
de preuve reste identique.

\textbf{Statut.} Assainissement (collecte verte rétablie) ; aucun axiome ajouté,
\texttt{theorie\_ensembles()==22}.


\subsection{§II.4.5 --- Commutativité de la réunion et de l'intersection binaires}
\textbf{Énoncé (Bourbaki).} Pour deux parties $A$ et $B$, la réunion et l'intersection sont
commutatives : $A \cup B = B \cup A$ et $A \cap B = B \cap A$.

\textbf{Formalisé.} Cibles : $A \cup B = B \cup A$ et $A \cap B = B \cap A$. \\
Module : \texttt{bourbaki/ensembles/familles/ii\_4\_reunion\_intersection\_familles/ii\_4\_1\_definitions\_algebre/ensembles\_binaire\_commut\_ii4.py}.

\textbf{Preuve (\Nstar).} Preuve point par point par extensionnalité. On instancie l'axiome de réunion
(resp. d'intersection) sur $z$ et l'on généralise (\texttt{generalisation}) : $A\cup B$ et $B\cup A$
sont caractérisés par la \emph{même} relation $R := (z\in A \vee z\in B)$, le côté $B\cup A$ étant
ramené à $R$ par \texttt{equivalence\_transitivite} avec la commutativité de $\vee$ (\texttt{comm\_ou}).
L'égalité s'obtient alors par \texttt{egalite\_par\_extension} (A1). Schéma identique pour $\cap$ avec
\texttt{comm\_et}. C'est la transposition à $\cup$/$\cap$ du schéma de $\{a,b\}=\{b,a\}$.

\textbf{Statut.} CLOS (0 hypothèse) pour les deux théorèmes ; \texttt{theorie\_ensembles()==22}.


\subsection{§II.4.5 (Cor. Prop. 6) --- Image directe d'une différence sous injection}
\textbf{Énoncé (Bourbaki).} Si $f$ est une injection de $A$ dans $B$, alors pour toute partie $X$ de $A$
l'image directe de la différence est la différence des images : $f\langle A \setminus X\rangle =
f\langle A\rangle \setminus f\langle X\rangle$.

\textbf{Formalisé.} Cible : $\mathrm{Injective}(f) \Rightarrow f\langle A \setminus X\rangle =
f\langle A\rangle \setminus f\langle X\rangle$. \\
Module : \texttt{bourbaki/ensembles/familles/ii\_4\_reunion\_intersection\_familles/ii\_4\_image\_famille/ensembles\_image\_difference\_injective.py}.

\textbf{Preuve (\Nstar).} Dual exact de l'image réciproque de la différence : les couples sont $(x,a)\in f$
(antécédent $x$, valeur $a$) et l'univalence devient l'injectivité. Preuve point par point sous
$\mathrm{Injective}(f)$ assumée (\texttt{assume}), via \texttt{membre\_image} (axiome d'image) et
\texttt{\_instance\_diff} (axiome de différence). L'inclusion $\supseteq$ est inconditionnelle. L'inclusion
$\subseteq$ utilise l'injectivité par le sous-lemme $(a\in f\langle X\rangle \Rightarrow x\in X)$ :
témoin $x'$, instanciation de l'injectivité sur $(x,a)$ et $(x',a)$ donnant $x=x'$, puis transport de
$x'\in X$ en $x\in X$ par Leibniz (\texttt{s6}). On combine \texttt{conjonction\_intro/elim},
\texttt{equivalence\_avant/arriere}, \texttt{contraposition}, \texttt{existe\_elimination} et les témoins
\texttt{s5} ; la double inclusion généralisée (\texttt{generalisation}) donne l'égalité par
\texttt{extensionnalite\_appliquee}. L'hypothèse $\mathrm{Injective}(f)$ est enfin déchargée par
\texttt{loi\_deduction}.

\textbf{Statut.} CLOS (0 hypothèse) : l'injectivité est portée dans la conclusion ($\mathrm{Injective}(f)
\Rightarrow \dots$) puis déchargée ; \texttt{theorie\_ensembles()==22}.


\subsection{§II.4.2 (Prop. 2) --- Associativité de la réunion d'une famille}
\textbf{Énoncé (Bourbaki).} Soit $(X_\iota)_{\iota\in L}$ une famille indexée et $L = \bigcup_\lambda J_\lambda$
une partition de l'ensemble d'indices. Alors la réunion peut s'associer selon les blocs :
$\bigcup_{\lambda\in L} X_\lambda = \bigcup_{\lambda\in L}\bigl(\bigcup_{\iota\in J_\lambda} X_\iota\bigr)$.
On en formalise la partie inconditionnelle (sans l'hypothèse $J_\lambda \neq \varnothing$).

\textbf{Formalisé.} Cible : $\bigcup_{\lambda\in L} X_\lambda = \bigcup_{\lambda\in L}\bigl(\bigcup_{\iota\in J_\lambda} X_\iota\bigr)$. \\
Module : \texttt{bourbaki/ensembles/familles/ii\_4\_reunion\_intersection\_familles/ii\_4\_1\_definitions\_algebre/ensembles\_familles\_assoc\_reunion.py}.

\textbf{Preuve (\Nstar).} L'hypothèse $L = \bigcup_\lambda J_\lambda$ est scindée fidèlement en deux clauses
assumées (\texttt{assume}) : \emph{(a) couverture} $(\forall\iota)(\iota\in L \Rightarrow (\exists\lambda)
(\lambda\in L \wedge \iota\in J_\lambda))$ et \emph{(b) domaine} $(\forall\lambda)(\forall\iota)((\lambda\in L
\wedge \iota\in J_\lambda) \Rightarrow \iota\in L)$. La famille interne $G_\lambda := \bigcup_{\iota\in J_\lambda}
X_\iota$ est définie par un terme (patron \texttt{famille\_reparam}), son axiome de valeur vivant dans une
théorie dédiée --- d'où aucune charge sur la théorie principale. La preuve est une double inclusion au point
$z$ : sens $\subseteq$ par la couverture (témoin $\lambda$, \texttt{s5} et \texttt{existe\_elimination}
imbriqués) ; sens $\supseteq$ après $\alpha$-renommage du $\exists$ externe (\texttt{alpha\_existe}, pour
éviter la capture du lieur interne par $J_\lambda$) et application du domaine. On y emploie \texttt{instancie},
\texttt{equivalence\_avant/arriere}, \texttt{conjonction\_intro/elim} et Leibniz (\texttt{s6}) pour
$G_\lambda = \bigcup_{\iota\in J_\lambda} X_\iota$. L'égalité finale provient de
\texttt{extensionnalite\_appliquee}.

\textbf{Statut.} CLOS sous hypothèses honnêtes $\{$\emph{couverture}~$(\forall\iota)(\iota\in L \Rightarrow
(\exists\lambda)(\lambda\in L \wedge \iota\in J_\lambda))$, \emph{domaine}~$(\forall\lambda)(\forall\iota)
((\lambda\in L \wedge \iota\in J_\lambda) \Rightarrow \iota\in L)\}$ ; la conclusion est l'égalité verbatim
(jamais prise en hypothèse) et n'est pas vacue ; \texttt{theorie\_ensembles()==22}.


\subsection{§II.4.2 --- Monotonie décroissante de l'intersection en l'indice}
\textbf{Énoncé (Bourbaki).} Si $J\subset I$, alors
$\bigcap_{\iota\in I}X_\iota \subset \bigcap_{\iota\in J}X_\iota$.

\textbf{Formalisé.} Cible : $(J\subset I) \Rightarrow (\bigcap_{\iota\in I}X_\iota \subset \bigcap_{\iota\in J}X_\iota)$. \\
Module : \texttt{.../ii\_4\_reunion\_intersection\_familles/ii\_4\_2\_proprietes/ensembles\_familles\_mono\_indice\_inter.py}.

\textbf{Preuve (\Nstar).} Dual universel ($\forall$, sans témoin) de
\texttt{reunion\_incluse\_sous\_indices} : on restreint le quantificateur via
$J\subset I$, \texttt{equivalence\_avant}/\texttt{arriere} de la caractérisation
de $\bigcap$, double \texttt{generalisation}.

\textbf{Statut.} CLOS ; \texttt{theorie\_ensembles()==22}.
